\documentclass[11pt,a4paper]{article}
\usepackage[margin=1in]{geometry}
\usepackage{hyperref}
\usepackage{enumitem}

\title{Strategic Management \\ Week 10 \\ Slide-by-Slide Notes (ELI5) \\ Unit 3}
\author{(Generated from \texttt{sm-week-10-slides-4.pdf})}
\date{}

\begin{document}
\maketitle

\section*{How to use this (ELI5)}
\begin{itemize}
  \item For every slide, write the ``Exam write'' line in your notebook.
  \item Then copy the ELI5 explanation in simple English.
  \item If you do not understand one slide, re-read only that part.
  \item Target: label + reason + exam sentence.
\end{itemize}

% ---------------------------
\section*{Slide 1 of 24 (Contents: Unit 3)}
\textbf{Key idea (from slide):} Unit 3 is Environment Analysis.\\
\textbf{ELI5:} You learn how outside world conditions your strategy.\\
\textbf{Exam write:} ``Unit 3 teaches how to diagnose opportunities/threats using macro and industry analysis.''

% ---------------------------
\section*{Slide 2 of 24 (Learning objectives)}
\textbf{Key idea (from slide):}
\begin{itemize}
  \item understand how environment conditions strategy and competitiveness
  \item analyze environmental factors beyond economics
  \item determine industry attractiveness (opportunities + threats)
\end{itemize}
\textbf{ELI5:} You need to understand the environment so you do not choose a strategy blindly.\\
\textbf{Exam write:} ``In answers, I will link environment diagnosis to strategy choices.''

% ---------------------------
\section*{Slide 3 of 24 (Business environment definition)}
\textbf{Key idea (from slide):}
\begin{itemize}
  \item firms operate inside political, social, cultural, legal systems
  \item management evaluates opportunities and threats
  \item environment = external characteristics outside firm control
  \item two levels: macro-environment and industry environment
\end{itemize}
\textbf{ELI5:} The environment is everything outside the firm that can help or hurt you.\\
\textbf{Exam write:} ``Environment analysis means finding outside factors the firm cannot fully control but that affect results.''

% ---------------------------
\section*{Slide 4 of 24 (Macro vs industry environments map)}
\textbf{Key idea (from slide):} Macro environment = political/social/tech/economic/legal/ecological. Industry environment = current competitors/customers/suppliers/potential competitors.\\
\textbf{ELI5:} Macro is the big background. Industry is who fights you and who buys from you.\\
\textbf{Exam write:} ``Split diagnosis into macro forces (PESTEL) and industry forces (Five Forces).''

% ---------------------------
\section*{Slide 5 of 24 (General environment analysis = PESTEL)}
\textbf{Key idea (from slide):} General environment analysis uses PESTEL.\\
\textbf{ELI5:} PESTEL is a checklist to find outside threats/opportunities.\\
\textbf{Exam write:} ``I will use PESTEL to translate macro factors into strategy-relevant opportunities and threats.''

% ---------------------------
\section*{Slide 6 of 24 (Exhibit 3.1 impact of Ukraine invasion)}
\textbf{Key idea (from slide):} Ukraine invasion impact is an example of how macro events change firms.\\
\textbf{ELI5:} Big political events can change costs, demand, and risk for many companies.\\
\textbf{Exam write:} ``In examples, I will show how a macro event changes costs/demand/risk.''

% ---------------------------
\section*{Slide 7 of 24 (Exhibit 3.2 impact of climate change)}
\textbf{Key idea (from slide):} Climate change impact is another macro example.\\
\textbf{ELI5:} Climate change affects industries through new risks and new opportunities.\\
\textbf{Exam write:} ``Climate change is diagnosed with PESTEL because it is political/economic/environment/technology at once.''

% ---------------------------
\section*{Slide 8 of 24 (Industry environment idea)}
\textbf{Key idea (from slide):}
\begin{itemize}
  \item Industrial Organization: industry structure affects strategy and profitability
  \item aim: discover ``industry attractiveness''
  \item use Five Forces analysis
  \item Five Forces = rivalry, entry threat, substitutes, buyer power, supplier power
  \item also defines what an industry is (product market + factors-of-production market)
\end{itemize}
\textbf{ELI5:} Industry analysis asks: ``How hard is it to make profit here?''\\
\textbf{Exam write:} ``Industry attractiveness is diagnosed with Five Forces, then translated into strategy.'

% ---------------------------
\section*{Slide 9 of 24 (Five Forces diagram)}
\textbf{Key idea (from slide):} Five Forces diagram shows the structure around rivals.\\
\textbf{ELI5:} Think of Five Forces as five ``pressure buttons'' around your firm.\\
\textbf{Exam write:} ``I will use the five pressures to explain why profit is easy or hard in an industry.''

% ---------------------------
\section*{Slide 10 of 24 (Degree of current competitive rivalry)}
\textbf{Key idea (from slide):}
\begin{itemize}
  \item rivalry among existing firms increases unattractiveness
  \item rivalry depends on: number/concentration (CR4), growth in demand, excess capacity + exit barriers, differentiation, cost conditions and large minimum effective size
 \end{itemize}
\textbf{ELI5:} If rivals fight hard, profits go down.\\
\textbf{Exam write:} ``To judge rivalry, I will use concentration, demand growth, exit barriers, differentiation, and minimum efficient scale.''

% ---------------------------
\section*{Slide 11 of 24 (Exhibit 3.3 current rivalry: concentration)}
\textbf{Key idea (from slide):} Exhibit about concentration (CR4 idea).\\
\textbf{ELI5:} Concentration tells you how many big players exist.\\
\textbf{Exam write:} ``Higher concentration can change rivalry intensity, so I will use CR4 evidence when possible.''

% ---------------------------
\section*{Slide 12 of 24 (Exhibit 3.4 current rivalry: concentration)}
\textbf{Key idea (from slide):} Another concentration exhibit example.\\
\textbf{ELI5:} More evidence to show how concentration supports your rivalry diagnosis.\\
\textbf{Exam write:} ``I will support rivalry claims with concentration-style evidence.''

% ---------------------------
\section*{Slide 13 of 24 (Threat of new entrants factors)}
\textbf{Key idea (from slide):}
\begin{itemize}
  \item if new competitors enter, rivalry increases and attractiveness drops
  \item factors: barriers to entry (administrative vs natural), government/legal barriers, capital requirements, economies of scale, absolute cost advantages or product differentiation, and access to distribution channels
  \item incumbents reaction: retaliation
\end{itemize}
\textbf{ELI5:} New entrants matter because they steal sales and force price pressure.\\
\textbf{Exam write:} ``I will evaluate entry threat using barriers (capital/scale/regulation) and incumbents' retaliation risk.''

% ---------------------------
\section*{Slide 14 of 24 (Exhibit 3.5 administrative barriers of entry)}
\textbf{Key idea (from slide):} Example of tariffs/regulation to protect industries and limit entry.\\
\textbf{ELI5:} Government rules can make entry harder (or easier).\\
\textbf{Exam write:} ``Administrative barriers can reduce entry threat by raising costs or requiring licenses.''

% ---------------------------
\section*{Slide 15 of 24 (Exhibit 3.6 natural barriers of entry)}
\textbf{Key idea (from slide):} Example of natural barriers via economies of scale.\\
\textbf{ELI5:} Sometimes big scale naturally blocks entry because new firms cannot compete on cost fast.\\
\textbf{Exam write:} ``Natural barriers like economies of scale can make entry unlikely and protect incumbents.''

% ---------------------------
\section*{Slide 16 of 24 (Threat of substitute products)}
\textbf{Key idea (from slide):}
\begin{itemize}
  \item substitutes meet the same needs with different technology
  \item substitutability depends on customers, perceived appreciation, and relative price
  \item more substitution threat means the sector is less attractive
  \item sources: disruptive technology, disruptive consumer behavior
\end{itemize}
\textbf{ELI5:} Substitutes are ``other ways'' to satisfy the same need, so they cap your ability to charge.\\
\textbf{Exam write:} ``Substitution risk is judged by customer acceptance and relative price of alternatives.''

% ---------------------------
\section*{Slide 17 of 24 (Illustration 3.1: explore new substitutes)}
\textbf{Key idea (from slide):} Task: explore up to 5 new substitute products and reasons (technology push and market/customer push).\\
\textbf{ELI5:} Practice brainstorming alternatives that could replace you.\\
\textbf{Exam write:} ``To answer substitutes, I will brainstorm alternatives by technology push and customer push.''

% ---------------------------
\section*{Slide 18 of 24 (Exhibit 3.7 substitutes)}
\textbf{Key idea (from slide):} Exhibit about substitute products.\\
\textbf{ELI5:} Examples help you see what ``substitute'' really means.\\
\textbf{Exam write:} ``Use examples to show the substitute threat matches need-level competition.''

% ---------------------------
\section*{Slide 19 of 24 (Bargaining power of buyers and suppliers)}
\textbf{Key idea (from slide):}
\begin{itemize}
  \item buyer power = customers' ability to negotiate trading conditions in their advantage
  \item factors: buyer concentration relative to industry, differentiation of traded products, information, threat of upstream integration
  \item supplier power logic mirrors buyers: if suppliers have power, industry attractiveness drops
end{itemize}
\textbf{ELI5:} If buyers have strong power, they push prices down. If suppliers have strong power, they push costs up.\\
\textbf{Exam write:} ``Buyer power and supplier power affect profit margins by shifting price-cost terms.''

% ---------------------------
\section*{Slide 20 of 24 (Importance of the environment: industry vs firm effect)}
\textbf{Key idea (from slide):} Industry vs firm effect explains how much outcomes come from environment vs from firm-specific factors (variance weights).\\
\textbf{ELI5:} Sometimes industry explains success, sometimes the firm explains success.\\
\textbf{Exam write:} ``In answers, I will separate environment impact from firm-specific capability impact.''

% ---------------------------
\section*{Slide 21 of 24 (Critics: are all Five Forces equally relevant?)}
\textbf{Key idea (from slide):}
\begin{itemize}
  \item critics ask: Are all forces equally relevant?
  \item role of complementary products
  \item public administrations
  \item industry segmentation and strategic groups
end{itemize}
\textbf{ELI5:} The model is a guide; sometimes one force matters more than others.\\
\textbf{Exam write:} ``I will discuss which Five Force is strongest, and mention complements/segmentation when relevant.''

% ---------------------------
\section*{Slide 22 of 24 (Other environmental issues: clusters)}
\textbf{Key idea (from slide):} Clusters/industrial districts and geographic agglomeration: talent attraction, entrepreneurship/startups, technological innovation, and access to resources.\\
\textbf{ELI5:} Clusters are like teams in one neighborhood that help each other grow.\\
\textbf{Exam write:} ``Clusters can raise opportunity by improving talent, innovation, and access to resources.''

% ---------------------------
\section*{Slide 23 of 24 (Clusters description: agents inside a cluster)}
\textbf{Key idea (from slide):} Agglomeration can produce extra opportunities for firms inside: companies devoted to an activity, suppliers and related industries, and university/research centres.\\
\textbf{ELI5:} Nearby companies and universities make you learn faster.\\
\textbf{Exam write:} ``In cluster analysis, identify firms, suppliers, related industries, and research institutions as sources of advantage.''

% ---------------------------
\section*{Slide 24 of 24 (Illustration 3.2: cluster task)}
\textbf{Key idea (from slide):} Task: identify a well-developed cluster and describe main features/agents and associated advantages.\\
\textbf{ELI5:} This is practice: pick a real cluster and list why it works.\\
\textbf{Exam write:} ``For a cluster question, describe agents (firms/suppliers/universities) and explain the advantage mechanism.''

\end{document}

